\begin{description}
	\item[尺度化とは何か] 心理学で行われるアンケート調査やその後の分析はどういう原理があって「心を測定した」といえるのか。その原理やモデルを理解して利用できるようになる。
	\item[多変量解析から何がわかるのか] 調査研究などで得られた多変量を分析することで何がわかるのか。あるいは何をしてわかったというのか。
	\item[多変量解析の基礎となる数式的原理] 多変量解析の背景にあるのは線形代数という数学であり，線形代数の基礎を学ぶことで多変量解析のメカニズムを統合的に理解できる。
	\item[データ生成メカニズム] データから情報を取り出す受け身的な分析ではなく，データに数字を与えたり，データが生まれてくるメカニズムをリバースエンジニアリングすることで，さらに積極的にデータ解析に立ち向かおう。
	\item[統計環境Rと確率的プログラミング言語Stanによる実践] 統計環境Rと確率的プログラミング言語Stanに習熟することで実際に計算し，確認しながら分析を進めることができる。
\end{description}
